\section{The prolate trials and their Rayleigh upper bounds}
\label{sec:prolate-trial}
This section defines the original trial in the spectral programme and
checks the upper estimate used in the Rayleigh criterion.
Andrews, Connes and La Gatta record a related upper estimate of Connes
and Kastoris for a Gevrey modification of the trial
\cite[Theorem~3.1, p.~14]{andrews-connes-lagatta2026}.
Their map first multiplies the raw prolate vector by an even cutoff in
an endpoint layer of width $\lambda^{-3}$, adds two fixed smooth bumps
to restore its value and integral at zero, applies the theta map, and
cuts off the lower logarithmic endpoint in a layer of width
$\lambda^{-4}$. They give a Rayleigh bound
$C(1-\mu_4)(\log\lambda)^3$ for that vector. In the notation below,
$1-\mu_4^2=(1+\mu_4)(1-\mu_4)$ exactly. We give the proof for the
original sharply truncated vector, whose exact definition is needed
in the Rayleigh criterion.
For raw variables use
$\mathcal Ff(v)=\int f(u)e^{-2\pi iuv}\,du$ and
$P_\lambda={\bf1}_{[-\lambda,\lambda]}$, with $L=\log\lambda$.
Let $\phi_{j,\lambda}$ be the real $L^2$-unit time-supported prolate
functions with Hermite index $j$, and write
\[
 P_\lambda\mathcal F\phi_{j,\lambda}
       =\mu_j\phi_{j,\lambda},\qquad
 \nu_j=\mu_j^2,\qquad d_j=\phi_{j,\lambda}(0).
\]
For $j=0,4$, the signs are chosen so that the limiting Hermite
functions have their usual real signs. For large $\lambda$ the two
$\mu_j$ are positive and $d_j\ne0$. Define
\[
 \psi_\lambda=-\frac34\left(
       \mu_4\frac{\phi_{0,\lambda}}{d_0}
       -\mu_0\frac{\phi_{4,\lambda}}{d_4}\right),\qquad
 k_\lambda(x)={\bf1}_{[-L,L]}(x)e^{x/2}
                  \sum_{n\ge1}\psi_\lambda(ne^x).
\]
The sum is finite for each $x$. The even normalized trial is
$v_L=(k_\lambda(x)+k_\lambda(-x))/(2\|k_\lambda^+\|_2)$,
where $k_\lambda^+$ denotes the even part.
This is the two-mode construction of
\cite[\S7, pp.~28--31]{ccm2025}, with its scalar fixed as follows.

The Hermite ratios at zero are
\[
 h_0(u)/h_0(0)=e^{-\pi u^2},\qquad
 h_4(u)/h_4(0)=
 \tfrac13(16\pi^2u^4-24\pi u^2+3)e^{-\pi u^2}.
\]
Their difference times $-3/4$ is
$h(u)=(4\pi^2u^4-6\pi u^2)e^{-\pi u^2}$.
Consequently
\[
 K(x)=e^{x/2}\sum_{n\ge1}h(ne^x),\qquad
 \widehat K(t)=\xi(1/2-it)=\Xi(t).
\]
For the transform identity, the Mellin integral of $h$ is
$\tfrac12s(s-1)\pi^{-s/2}\Gamma(s/2)$; multiply by $\zeta(s)$
and continue the absolutely convergent theta integral.
Poisson summation, $\mathcal Fh=h$, $h(0)=0$ and $\int h=0$ show
$K(-x)=K(x)$. Its displayed series gives decay faster than every
exponential at positive infinity, and hence at both ends.
Thus its Fourier integral is entire and $K\ne0$.

The fixed-index Hermite comparison in
\cite[Lemma~7.2, p.~29]{ccm2025} gives
$\sup_{|u|\le\lambda}|\psi_\lambda(u)-h(u)|\ll\lambda^{-2}$.
For completeness, the time-supported unit normalizations differ from
their Hermite limits in $L^2$ by $O(\lambda^{-3/2})$. Since the modes
$0,4$ are Fourier invariant in the limit, their leakage tends to zero;
this also gives $\mu_j\to1$ and the stated ratio comparison at zero.
Summing the uniform error over $n\le\lambda e^{-x}$ gives
$|k_\lambda(x)-K(x)|\ll\lambda^{-1}e^{-x/2}$ on $[-L,L]$, apart
from a Gaussian tail whose $L^2$ norm is exponentially small.
Integration yields
\begin{equation}
 \|k_\lambda-K\|_2=O(\lambda^{-1/2}),\qquad
 v_L\longrightarrow K/\|K\|_2.
 \label{eq:prolate-trial-limit}
\end{equation}
The raw mean is exactly zero: evaluation of the finite Fourier equation
at zero gives $\int_0^\infty\phi_{j,\lambda}=\mu_jd_j/2$, and these
terms cancel in $\psi_\lambda$.

\begin{proposition}[Leakage controls the original trial]
\label{prop:prolate-upper}
Put $\delta_\lambda=1-\nu_4$. For every $\varepsilon>0$,
\[
 q(k_\lambda)\ll_\varepsilon\lambda^\varepsilon\delta_\lambda,
 \qquad q(v_{\log\lambda})
                  \ll_\varepsilon\lambda^\varepsilon\delta_\lambda.
\]
The classical fixed-index prolate asymptotic gives
\[
 \delta_\lambda\sim
  \frac{2^{15}\sqrt2\pi^5}{3}\lambda^9e^{-4\pi\lambda^2},
\]
and consequently
$a_1(\log\lambda)\ll_\varepsilon
\lambda^{9+\varepsilon}e^{-4\pi\lambda^2}$.
\end{proposition}

\begin{proof}
We first prove the exterior estimates needed for the logarithmic tail.
Fix an even index $j$ and write $y=\mathcal F\phi_{j,\lambda}$,
$\delta_j=1-\nu_j$. Its differential eigenvalue $\chi_j$ obeys
$0\le\chi_j\ll_j\lambda^2$: a fixed-dimensional space of compact
smooth test functions proves the upper bound by min--max in the
prolate differential form. For $x>\lambda$, the Fourier continuation
satisfies
\[
 (py')'+Vy=0,\qquad
 p=x^2-\lambda^2,\qquad V=4\pi^2\lambda^2x^2-\chi_j.
\]
To justify the equation at the endpoint, extend the time-supported
function by zero. Multiplication of its distributional derivative by
$\lambda^2-x^2$ removes the endpoint delta terms; the remaining flux
also vanishes at both endpoints. Fourier transformation therefore gives
the equation globally, and $y$ is entire.
For $E=p|y'|^2+V|y|^2$, direct differentiation gives
\[
 (E/V)'=-\left(p'/V+pV'/V^2\right)|y'|^2\le0.
\]
Thus $|y(x)|\le|y(\lambda)|$. Integrating the equation from the
endpoint gives $|y'(x)|\ll\lambda^3|y(\lambda)|$ for
$\lambda\le x\le2\lambda$.

Here is an endpoint-to-leakage estimate with no differentiation of an
asymptotic formula. Put $c=2\pi\lambda^2$, $x=\lambda\cosh t$, and
$Z(t)=\sqrt{\sinh t}\,y(\lambda\cosh t)$. Then
\[
 Z''+\Omega Z=0,\qquad
 \Omega=c^2\cosh^2t-\chi_j-\tfrac14+(4\sinh^2t)^{-1}.
\]
For a fixed small $a>0$, set $t_0=a/c$ and $t_1=\operatorname{arcosh}2$.
The preceding derivative bound gives
$|y(\lambda\cosh t_0)|\ge|y(\lambda)|/2$.
On $[t_0,t_1]$, $\Omega$ lies between positive multiples of $c^2$ and
has one minimum, because
$\Omega'=\cosh t\,(4c^2\sinh^4t-1)/(2\sinh^3t)$.
The energy $\mathcal E=|Z'|^2+\Omega|Z|^2$ satisfies
$\mathcal E'=\Omega'|Z|^2$. On the decreasing portion use
$\mathcal E'\ge(\Omega'/\Omega)\mathcal E$; on the increasing
portion use $\mathcal E'\ge0$. The bounded ratio of the extrema gives
$\mathcal E(t_1)\gg c|y(\lambda)|^2$.

For $U(x)=xy(x)$ one obtains
\[
 U''+bU'+k^2U=0,\quad
 b=\frac{2\lambda^2}{x(x^2-\lambda^2)},\quad
 k^2=4\pi^2\lambda^2+
 \frac{4\pi^2\lambda^4-\chi_j-2\lambda^2/x^2}{x^2-\lambda^2}.
\]
On $x\ge2\lambda$, $k^2\ge4\pi^2\lambda^2$ and $(k^2)'<0$.
Hence $\mathcal E_U=|U'|^2+k^2|U|^2$ decreases.
At $2\lambda$ it is comparable to $\mathcal E(t_1)$, both being
comparable to
$\lambda^2|y'(2\lambda)|^2+\lambda^4|y(2\lambda)|^2$.
On $[2\lambda,3\lambda]$ the differential inequality
$\mathcal E_U'\ge(-2b+(k^2)'/k^2)\mathcal E_U$ shows that its
values are comparable. Choose a smooth cutoff $\theta$ in that
interval equal to one on its middle half, with
$|\theta'|\ll\lambda^{-1}$. Multiplication of the equation by
$\theta^2\overline U$ and integration by parts give
\[
 \int\theta^2|U'|^2
 =\int\theta^2k^2|U|^2+
  \operatorname{Re}\int(\theta^2b-2\theta\theta')U'\overline U.
\]
Young's inequality absorbs half the derivative integral and bounds
the rest by $O(\lambda^2\int|U|^2)$. Thus
$\lambda\mathcal E_U(2\lambda)
\ll\lambda^2\int_{2\lambda}^{3\lambda}|U|^2$, which yields
\[
 \delta_j\gg\lambda^{-3}\mathcal E_U(2\lambda)
                  \gg\lambda^{-1}|y(\lambda)|^2.
\]
The upper energy estimate at $2\lambda$, followed by its decrease,
also gives
\[
 |y(x)|\ll\lambda|y(\lambda)|/x,\qquad
 |y'(x)|\ll\lambda^4|y(\lambda)|/x\quad(x\ge\lambda).
\]
The interval $[\lambda,2\lambda]$ is covered by the earlier bounds.

Define the sharply cut logarithmic Fourier tail
$w_j(t)={\bf1}_{t>0}\sqrt\lambda e^{t/2}y(\lambda e^t)$.
Exactly $\|w_j\|_2^2=\delta_j/2$, while the exterior bounds give
\[
 |w_j(t)|\le A e^{-t/2},\quad |w_j'(t)|\le B e^{t/2},\quad
 A\ll\lambda\sqrt{\delta_j},\quad
 B\ll\lambda^5\sqrt{\delta_j}.
\]
For every fixed integer $m\ge0$ and $0<s<1/2$, a translation-difference
estimate gives
$\|t^mw_j\|_{H^s}^2\ll_{j,m,s}A^{2-2s}B^{2s}$.
Here are the details including the jump at zero. For a translation
$0<r<A/B$, split the positive half-line at
$R=\log(A/(Br))$. On $[0,R]$ integrate the derivative bound; on
$[R,\infty)$ use the amplitude bound. The squared difference,
including the interval crossing zero, is at most
$C_mABr(1+|\log(A/(Br))|)^{2m}$.
For $r\ge A/B$ it is at most $C_mA^2$.
Integration against $r^{-1-2s}\,dr$ gives the claim, since
$\int_0^1r^{-2s}(1+|\log r|)^{2m}\,dr<\infty$.
In particular $\|t^mw_j\|_{H^{1/4}}\ll\lambda^2\sqrt{\delta_j}$.

The exact $L^2$ norm improves the size bound. Split at $4\log\lambda$;
use the exact norm before the split and the amplitude bound afterwards.
This gives
$\|t^mw_j\|_2\ll\sqrt{\delta_j}(1+\log\lambda)^m$.
Fourier H\"older interpolation therefore proves, for $0<s<1/4$,
\begin{equation}
 \|t^mw_j\|_{H^s}
 \ll_{j,m,s}\sqrt{\delta_j}\lambda^{8s}
                     (1+\log\lambda)^{m(1-4s)}.
 \label{eq:prolate-log-tail}
\end{equation}
The two coefficients of $\psi_\lambda$ are bounded, and
$\delta_0\le\delta_4$. Thus the same estimate, with
$\sqrt{\delta_\lambda}$, applies to the logarithmic tail $w$ of
$\mathcal F\psi_\lambda$ and its first moment.

For clarity we describe the theta operation on this tail, including
the change of variables at the lower endpoint. For a raw function on
the positive half-line put
$\mathcal Uf(t)=e^{t/2}f(e^t)$. This is an isometry from
$L^2(0,\infty;du)$ to $L^2(\mathbb R;dt)$. In logarithmic coordinates
let $\mathcal E$ be the closed Fourier multiplier
$\zeta(1/2-i\tau)$. For compactly supported functions of bounded
variation, $(\mathcal E\mathcal Uf)(\log u)$ equals
\[
 Ef(u)=\sqrt u\sum_{n\ge1}f(nu)-u^{-1/2}\int_0^\infty f(v)\,dv.
\]
This follows by integrating from $u=\epsilon$ and using
$\sum_{n\le Y}n^{-s}=Y^{1-s}/(1-s)+\zeta(s)+O_s(Y^{-\Re s})$;
the first term cancels the subtracted integral, and the remaining
error tends to zero for $0<\Re s<1$.
The bound of Chandee and Soundararajan
\cite[Theorem~1.1, p.~1]{chandee-soundararajan2009} gives, under the
Riemann hypothesis,
$|\zeta(1/2+i\tau)|+|\zeta'(1/2+i\tau)|
\ll_b(1+|\tau|)^b$ for every $b>0$.
Here the derivative is obtained as follows. Their bound on the critical
line is smaller than every positive power of $1+|\tau|$. On the line
$\Re s=2$, the Dirichlet series is bounded. The three-lines theorem,
applied after cancelling the pole with $(s-1)/(s+1)$, gives the same
subpower bound throughout $1/2\le\Re s\le2$. The functional equation
then bounds $\zeta$ in $1/2-d\le\Re s\le1/2$ by an additional factor
$O_d((1+|\tau|)^d)$. Choose a fixed $d<\min(b/2,1/4)$, use exponent $b/2$
in the preceding subpower bounds, and apply Cauchy's formula on circles
of radius $d/2$ about $1/2+i\tau$. The bounded range of $\tau$ is
handled by compactness away from the pole. This proves the displayed
derivative estimate. Consequently, for arbitrarily small positive $s',b$,
\[
 \|\mathcal Ew\|_{H^{s'}}+\|t\mathcal Ew\|_{H^{s'}}
 \ll_{s',b}\|w\|_{H^{s'+b}}+\|tw\|_{H^{s'+b}}.
\]
The second term follows by differentiating the Fourier multiplier.
Write $\varphi=\mathcal F\psi_\lambda$, $g=P_\lambda\varphi$ and
$\eta=(1-P_\lambda)\varphi$. If
$\mathcal Mf(s)=\int_0^\infty f(u)u^{s-1}\,du$, the cosine Mellin
identity gives
\[
 \mathcal M\varphi(s)
 =2(2\pi)^{-s}\Gamma(s)\cos(\pi s/2)
                         \mathcal M\psi_\lambda(1-s).
\]
It is justified by inserting an Abel factor in the cosine integral
and then removing it: $\psi_\lambda$ is bounded and compactly
supported, and its bounded-variation Fourier transform is $O(1/u)$,
so the final integrals converge for $0<\Re s<1$.
The zeta functional equation says that the displayed gamma factor
times $\zeta(s)$ equals $\zeta(1-s)$. On $s=1/2-i\tau$, the theta
multiplier identity therefore gives, by Plancherel,
\[
 E\psi_\lambda(e^x)=E\varphi(e^{-x}).
\]
Here $E\varphi$ means the closed multiplier applied to
$\mathcal U\varphi$; no pointwise sum over the noncompact Fourier
tail is required. Both sides belong to $L^2(dx)$: the compact raw
terms have logarithmic Sobolev regularity of every order below
$1/2$, and the tail has the regularity already proved above.

Because $\int_0^\infty\psi_\lambda=0$, the full theta function
$F_\lambda(x):=E\psi_\lambda(e^x)$ equals the uncorrected sum in
the definition of $k_\lambda$; it vanishes for $x>L$.
For $y>\lambda$, the finite sum defining $Eg(y)$ is zero, and hence
\[
 E\varphi(y)=E\eta(y)-c_Qy^{-1/2},\qquad
 c_Q=\int_0^\infty g(u)\,du
       =-\tfrac38\mu_0\mu_4(\mu_0-\mu_4).
\]
For the last equality, on the time interval
$g=-\tfrac34\mu_0\mu_4(\phi_{0,\lambda}/d_0-
\phi_{4,\lambda}/d_4)$; the two integrals are $\mu_0/2$ and
$\mu_4/2$. Thus $|c_Q|\ll\delta_\lambda$.

The logarithmic tail previously denoted by $w$ is exactly
$\mathcal U\eta_\lambda$, where
$\eta_\lambda(u)=\sqrt\lambda\,\eta(\lambda u)$ on $u>0$.
The theta operation satisfies the exact dilation identity
$E\eta_\lambda(y)=E\eta(\lambda y)$, first on compact raw inputs
by substitution in the sum and integral, and then on its closed
multiplier domain. For the lower tail set $x=-L-t$, so that
$t>0$ and $e^{-x}=\lambda e^t$. The preceding identities give
\begin{equation}
 R(t):={\bf1}_{t>0}F_\lambda(-L-t)
 = {\bf1}_{t>0}\left(\mathcal Ew(t)
                    -c_Q\lambda^{-1/2}e^{-t/2}\right).
 \label{eq:prolate-lower-tail}
\end{equation}

For $0<s<s'<1/2$, sharp restriction obeys
$\|{\bf1}_{t>0}f\|_{H^s}\ll_{s,s'}\|f\|_{H^{s'}}$.
The extra Gagliardo term is a constant times
$\int_0^\infty|f(t)|^2t^{-2s}\,dt$; on $(0,1)$ apply Sobolev
embedding into $L^{2/(1-2s')}$ and H\"older, and elsewhere use $L^2$.
Combining this with \eqref{eq:prolate-log-tail}, with all indices
sufficiently small, proves
\[
 \|R\|_{H^s}+\|tR\|_{H^s}
       \ll_\sigma\sqrt{\delta_\lambda}\lambda^{8\sigma}
                          (1+\log\lambda).
\]

Finally the local zero count, with multiplicities, is
$\sum_{\gamma\in[j,j+1]}m_\gamma\ll\log(2+|j|)$.
For $f,tf\in H^s$ the one-dimensional $H^1$ pointwise bound for
$\widehat f$ on each unit interval and Plancherel give
\[
 \sum_\gamma m_\gamma|\widehat f(\gamma)|^2
      \ll_s\|f\|_{H^s}^2+\|tf\|_{H^s}^2.
\]
Indeed bound the supremum of $|\widehat f|^2$ by the integrals of
$|\widehat f|^2+|\widehat f'|^2$ on a slightly larger interval,
multiply by its zero count, and use
$\log(2+|\tau|)\ll_s(1+\tau^2)^s$.
The Fourier transform of $F_\lambda$ has factor
$\zeta(1/2-i\tau)$ and hence vanishes at every ordinate. Its upper
tail is zero. The bounds for $R,tR$ make $R$ integrable by
Cauchy--Schwarz, so $F_\lambda$ is integrable as well. Both sides of
the Fourier multiplier identity are continuous on the real line;
their $L^2$ equality therefore holds at each ordinate.
The exact substitution $x=-L-t$ in its lower tail gives
\[
 \widehat{k_\lambda}(\gamma)
       =-e^{i\gamma L}\widehat R(-\gamma).
\]
Reflection preserves the signed sampling sum, so
$q(k_\lambda)=\sum_\gamma m_\gamma|\widehat R(\gamma)|^2$.
This yields
\[
 q(k_\lambda)\ll_\sigma
       \lambda^{16\sigma}(1+\log\lambda)^2\delta_\lambda.
\]
Choose $16\sigma<\varepsilon$ and absorb the logarithm.
Reflection is a $q$-isometry, so taking the even part cannot increase
$q$, and its norm tends to $\|K\|_2>0$ by
\eqref{eq:prolate-trial-limit}. This proves the two Rayleigh bounds.

To state their exponential scale, Fuchs's fixed-index theorem
\cite[Theorem~1, p.~319]{fuchs1964} has parameter $c=2\pi\lambda^2$
and concentration eigenvalue $\nu_j$:
$1-\nu_j\sim(4\sqrt\pi\,8^j/j!)c^{j+1/2}e^{-2c}$.
Substitution of $j=4$ gives the constant in the proposition.
Min--max and the nonzero trial norm prove the asserted bound for $a_1$.
\end{proof}

\begin{proposition}[An odd prolate competitor]
\label{prop:odd-prolate-upper}
For the prolate modes of indices $2,6$, whose finite Fourier
eigenvalues are negative, define
\[
 \psi_\lambda^{\rm o}
   =\frac{15}{8}\left(\mu_6\frac{\phi_{2,\lambda}}{d_2}
                  -\mu_2\frac{\phi_{6,\lambda}}{d_6}\right),
 \qquad
 k_\lambda^{\rm o}(x)={\bf1}_{[-L,L]}(x)e^{x/2}
                   \sum_{n\ge1}\psi_\lambda^{\rm o}(ne^x).
\]
Its odd part has nonzero norm for sufficiently large $\lambda$.
Writing $v_L^{\rm o}$ for that part divided by its $L^2$ norm, one has
\[
 \|k_\lambda^{\rm o}-K'\|_2=O(\lambda^{-1/2}),\qquad
 q(v_L^{\rm o})\ll_\varepsilon\lambda^\varepsilon(1-\mu_6^2).
\]
In particular,
\begin{equation}
 a_2(\log\lambda)\ll_\varepsilon
                  \lambda^{13+\varepsilon}e^{-4\pi\lambda^2}
 \quad(\varepsilon>0).
 \label{eq:second-eigenvalue-upper}
\end{equation}
\end{proposition}

\begin{proof}
The fixed-index Hermite limit in \cite[equation~(7.10), p.~30]{ccm2025}
applies to indices $2,6$ and gives the ratios
\[
 \frac{h_2(u)}{h_2(0)}=(1-4\pi u^2)e^{-\pi u^2},\qquad
 \frac{h_6(u)}{h_6(0)}=
 \left(1-12\pi u^2+16\pi^2u^4-\frac{64}{15}\pi^3u^6\right)
 e^{-\pi u^2}.
\]
Since $\mu_2,\mu_6\to-1$, their displayed combination tends to
\[
 \frac{15}{8}\left(\frac{h_6}{h_6(0)}-\frac{h_2}{h_2(0)}\right)
  =(-15\pi u^2+30\pi^2u^4-8\pi^3u^6)e^{-\pi u^2}
  =\left(u\frac{d}{du}+\frac12\right)h(u).
\]
Differentiating the absolutely convergent theta series for $K$ shows
that the theta image of this last function is exactly $K'$.
The same uniform Hermite comparison and summation used in
\eqref{eq:prolate-trial-limit} give the asserted $L^2$ error.
The function $K'$ is odd and nonzero, because $K$ is even, nonconstant,
and tends to zero at infinity.  This proves the norm assertion.

The mean of $\psi_\lambda^{\rm o}$ is exactly zero by
$\int_0^\infty\phi_j=\mu_jd_j/2$.  Every exterior estimate in the
preceding proof was established for an arbitrary fixed even index.
It therefore applies to the two bounded coefficients of the modes
$2,6$, with leakage $\delta^{\rm o}=1-\mu_6^2$.
There is one scalar term in the Poisson tail to check: for
$g=P_\lambda\mathcal F\psi_\lambda^{\rm o}$,
\[
 \int_0^\infty g(u)\,du
   =\frac{15}{16}\mu_2\mu_6(\mu_2-\mu_6)
   =O(\delta^{\rm o}).
\]
Both eigenvalues approach $-1$, and
$|\mu_j+1|=(1-\mu_j^2)/(1-\mu_j)$, which proves this bound with
the sign and denominator explicit.  The exact reflected-tail
identity, Sobolev estimate, zeta multiplier, and sampling argument
from the preceding proof now give
$q(k_\lambda^{\rm o})\ll_\varepsilon
\lambda^\varepsilon\delta^{\rm o}$.  Taking the odd part decreases
the form and its normalization stays bounded away from zero.

The even trial and this odd trial are orthogonal both in $L^2$ and
for $q$: pair the two signed ordinates in the sampling formula for
the latter assertion.  The largest Rayleigh quotient on their
two-dimensional span is consequently the larger of their two
quotients.  Min--max and the fixed-index Fuchs formula give
\[
 a_2(\log\lambda)\ll_\varepsilon
       \lambda^\varepsilon(1-\mu_6^2),\qquad
 1-\mu_6^2\sim\frac{2^{22}\sqrt2\pi^7}{45}
                          \lambda^{13}e^{-4\pi\lambda^2}.
\]
This proves \eqref{eq:second-eigenvalue-upper}.
\end{proof}

The baseline power in the Rayleigh criterion is therefore $9$.
Precisely, the change of parameter between an absolute power $\beta$
and a gain $\alpha$ above the trial scale is $\beta=9+\alpha$:
\[
 \lambda^\beta e^{-4\pi\lambda^2}
  =\lambda^\alpha\bigl(\lambda^9e^{-4\pi\lambda^2}\bigr).
\]
Since $(1-\mu_4^2)/(\lambda^9e^{-4\pi\lambda^2})$ tends to a
positive constant, positive-power improvement over that leakage is
exactly a positive gain $\alpha$.  Equation
\eqref{eq:second-eigenvalue-upper} rules out every eventual lower
bound $a_2(\log\lambda)\ge c\lambda^\beta e^{-4\pi\lambda^2}$
with $\beta>13$: choose $0<\varepsilon<\beta-13$ and divide the two
bounds.  Thus the positive-gain range still open by this comparison
is $9<\beta\le13$, equivalently $0<\alpha\le4$.
